

\subsection{Les niveaux d'isolation SQL}

\begin{frame}{Ce qu'on doit savoir }

\alert{Il existe plusieurs niveaux d'isolation, du plus permissif au plus strict}.

\vfill

Plus le niveau est permissif, plus l'exécution est fluide, plus les anomalies sont possibles.

\vfill

Plus le niveau est strict, plus l'exécution risque de rencontrer des blocages, moins
les anomalies sont possibles.

\vfill

\alert{Le niveau d'isolation par défaut n'est jamais le plus strict}. Car

\begin{redbullet}
  \item inutile dans la grande majorité des cas;
  \item provoque rejets et blocages difficiles à expliquer à l'utilisateur.  
\end{redbullet}

\vfill

\begin{blocgauche}
\alert{Quand l'isolation totale est nécessaire, il faut l'indiquer explicitement.}  
\end{blocgauche}
\end{frame}

\begin{frame}{Niveaux d'isolation SQL}

Définis en fonction des anomalies: lectures sales,  lectures non répétables, tuples fantômes.

\begin{redbullet}
  \item \texttt{read uncommitted}: tout est permis, toutes les anomalies sont possibles.
  
  \item \texttt{read committed}: lectures sales non permises.
  
  \item \texttt{repeatable read}: seuls les tuples fantômes sont permis.
  
  \item \texttt{serializable}: isolation totale, aucune anomalie, interblocages possibles.
\end{redbullet}

Niveau par défaut: \texttt{read committed} (Oracle) ou \texttt{repeatable read} (MySQL, PostgreSQL).
 
 \vfill
 
 
\begin{blocgauche}
\alert{Il faut se mettre en mode \texttt{serializable} pour les processus transactionnels.}
\end{blocgauche}

\end{frame}

\begin{frame}{Le mode \texttt{read committed}}

\alert{Retenir}: une requête accède à l'état de la base \alert{au moment où la
requête est exécutée}.

\vfill

Pas de lecture sale, car donnée en cours de modification ne fait pas partie de l'état de la base.

\vfill

Assez fluide, mais autorise beaucoup d'anomalies (mises à jour perdues, contrôle).
 
\vfill


\begin{blocgauche}
\alert{Démonstration}: on peut lire deux fois le même tuple dans une transaction et obtenir
des résultats différents.
\end{blocgauche}


\end{frame}


\begin{frame}{Le mode \texttt{repeatable read}}

\alert{Retenir}: une requête accède à l'état de la base \alert{au moment où la
transaction a débuté}.

\vfill

Pas de lecture sale, pas de lecture non répétable: les requêtes accèdent toujours
au même état de la base. 

\vfill

Autorise les mises à jour perdues.
 
\vfill



\begin{blocgauche}
\alert{Démonstration}: on peut lire $n$ fois le même tuple dans une transaction et obtenir
toujours le même résultat.
\end{blocgauche}


\end{frame}


\begin{frame}{Est-ce suffisant?}

\alert{Non}, car des exécutions non sérialisables restent quand même possibles.


\vfill


\alert{Démonstration}: On prend notre base et on déroule l'exécution des mises à jour perdues.

\end{frame}


\begin{frame}[fragile]{Le mode \texttt{serializable}}


\alert{Retenir}: garantit l'isolation totale, et donc la cohérence de la base.

\vfill

\alert{mais ...}: risque non négligeable de \alert{rejet} de l'une des transactions.


\vfill

On se place dans ce mode avec la commande suivante, au début du code de la transaction.

\begin{minted}[frame=leftline,
               baselinestretch=.8,
               fontsize=\small,
               framesep=2mm]{sql}
  set transaction isolation level serializable;
\end{minted}

\vfill

Démonstration...

\end{frame}

\begin{frame}{La clause \texttt{for update}}

\alert{Fait}: le problème vient du fait qu'on \alert{lit} une donnée pour la \alert{modifier} ensuite.

\begin{redbullet}
  \item Le système ne peut pas le deviner!
  \item Il est obligé de tracer toutes les lectures (coût élevé)
  \item Il pose des verrous faibles $\rightarrow$ rejet quand les choses tournent mal
\end{redbullet}

\vfill


La clause \texttt{for update}
\begin{redbullet}
  \item Le programmeur déclare qu'une lecture va être suivie d'une mise à jour
   \item Le système pose un verrou fort pour celles-là, pas de verrou pour les autres
\end{redbullet}

\vfill

\begin{blocgauche}
\alert{Permet de monter le niveau de blocage uniquement quand c'est nécessaire} ...
mais repose sur le facteur humain, pas assez fiable.
\end{blocgauche}
\end{frame}

\begin{frame}{À retenir}

Savoir repérer les transactions dans une application. Elle doivent respecter la 
\alert{cohérence applicative} (le système ne peut pas la deviner pour vous).

\vfill 

Dans ce cas, appliquer le mode sérialisable.

\vfill

La clause \texttt{for update}  est en théorie meilleure, mais repose 
sur le facteur humain: probablement à éviter.

\vfill

Savoir qu'en mode sérialisable on risque des rejets.

\vfill

Comprendre les risques dans les autres modes.

\end{frame}